% cs250preamble.tex --- shared preamble for CS 250 course notes
%
% This file is \input by main.tex and by each individual module
% file so that each module can still compile standalone. Any
% package or custom-environment change should be made here, not
% in the individual module files.

\usepackage[T1]{fontenc}
\usepackage[margin=1in]{geometry}
\usepackage{amsmath, amssymb, amsthm}
\usepackage{enumitem}
\usepackage{hyperref}
\usepackage{booktabs}
\usepackage{listings}
\usepackage{xcolor}
\usepackage{tikz}
\usetikzlibrary{shapes.geometric, shapes.gates.logic.US, shapes.gates.logic.IEC, arrows.meta, positioning, calc, trees}
\usepackage{bussproofs}
\usepackage{mdframed}

\lstset{
  basicstyle=\ttfamily\small,
  frame=single,
  breaklines=true,
  columns=fullflexible,
  mathescape=true
}

\theoremstyle{definition}
\newtheorem{theorem}{Theorem}[section]
\newtheorem{definition}[theorem]{Definition}
\newtheorem{example}[theorem]{Example}

\hypersetup{
  colorlinks=true,
  linkcolor=blue!60!black,
  urlcolor=blue!60!black
}

% Key-takeaway box: used at the end of major sections and at the end of
% each module to summarise the main points. Expects \item bullets inside.
\newmdenv[
  linewidth=0.8pt,
  linecolor=blue!40!black,
  backgroundcolor=blue!3,
  skipabove=8pt,
  skipbelow=8pt,
  innertopmargin=7pt,
  innerbottommargin=7pt,
  innerleftmargin=9pt,
  innerrightmargin=9pt
]{keytakeawaybox}
\newenvironment{keytakeaway}{%
  \begin{keytakeawaybox}%
  \noindent\textbf{Key takeaways.}%
  \begin{itemize}[leftmargin=1.3em, topsep=2pt, itemsep=1pt, label=\textbullet]%
}{%
  \end{itemize}%
  \end{keytakeawaybox}%
}

% Summary/looking-ahead environment: used once at the end of each module
% to wrap up what the module built and point forward.
\newmdenv[
  linewidth=0.8pt,
  linecolor=gray!60,
  backgroundcolor=gray!5,
  skipabove=8pt,
  skipbelow=8pt,
  innertopmargin=7pt,
  innerbottommargin=7pt,
  innerleftmargin=9pt,
  innerrightmargin=9pt
]{summarybox}

% Sidebar environment: used for tangential asides---historical notes,
% pointers to other areas of math/CS, cross-paradigm remarks---that
% enrich the main narrative without being essential to it. Takes one
% argument: the sidebar's title.
\newmdenv[
  linewidth=0.6pt,
  linecolor=gray!60,
  backgroundcolor=gray!4,
  skipabove=8pt,
  skipbelow=8pt,
  innertopmargin=6pt,
  innerbottommargin=6pt,
  innerleftmargin=9pt,
  innerrightmargin=9pt
]{sidebarbox}
\newenvironment{sidebar}[1]{%
  \begin{sidebarbox}%
  \noindent\textbf{Sidebar: #1.}\par\medskip%
}{%
  \end{sidebarbox}%
}

% Reading-map environment: used at the top of each numbered module to
% indicate Epp coverage, prerequisites, relevant intermezzi, and
% suggested practice problems. Expects four \rmentry{label}{body}
% items inside (or arbitrary paragraphs).
\newmdenv[
  linewidth=0.8pt,
  linecolor=black!55,
  backgroundcolor=yellow!4,
  skipabove=10pt,
  skipbelow=14pt,
  innertopmargin=9pt,
  innerbottommargin=9pt,
  innerleftmargin=11pt,
  innerrightmargin=11pt
]{readingmapbox}
\newcommand{\rmentry}[2]{\par\noindent\textbf{#1}\hspace{0.4em}#2\par}
\newenvironment{readingmap}{%
  \begin{readingmapbox}%
  \noindent{\bfseries\large Reading map}%
  \vspace{2pt}%
  \setlength{\parskip}{4pt plus 1pt}%
}{%
  \end{readingmapbox}%
}
